\chapter{Riproducibilità e inventario degli artefatti}
\label{app:reproducibility}

Questa appendice descrive come il repository GraphMate supporta la riproducibilità scientifica.  Evita deliberatamente di ripetere ogni iperparametro, feature del grafo e regola di benchmark.  Tali dettagli appartengono alle Appendici~\ref{app:hyperparameters}--\ref{app:llm-protocol}.  L'enfasi qui è posta sui confini degli artefatti, sui file destinati alla distribuzione pubblica e sulla differenza tra ispezione, valutazione di checkpoint congelati e riaddestramento.

\section{Classi di artefatti}
\label{app:reproducibility:artifacts}

GraphMate separa configurazioni, checkpoint, risorse, dataset generati e uscite di analisi.  Le configurazioni ufficiali definiscono architettura prevista e politica di addestramento.  I checkpoint congelati memorizzano i pesi addestrati selezionati e sono identificati da hash SHA-256.  Le risorse includono file ausiliari stabili, come il vocabolario delle mosse.  Gli artefatti contengono report di addestramento, uscite di valutazione e riepiloghi dell'analisi finale.  La Tabella~\ref{tab:artifact-inventory} sintetizza le principali classi di artefatti destinate alla riproducibilità.

\begin{table}[H]
\centering
\small
\caption{Classi di artefatti del repository rilevanti per la riproducibilità.}
\label{tab:artifact-inventory}
\begin{tabular}{p{0.22\linewidth}p{0.30\linewidth}p{0.34\linewidth}}
\toprule
Classe di artefatto & Esempi & Ruolo nella riproducibilità \\
\midrule
Configurazioni ufficiali & \path{configs/model_*.json} & Architettura, root dei dati, ottimizzatore, scheduler e politica di addestramento. \\
Checkpoint congelati & \path{checkpoints/model_*/best.pt} & Pesi pretrained identificati da SHA-256. \\
Risorse & JSON del move encoder & Mappature ausiliarie stabili, come il vocabolario delle mosse costruito sull'addestramento. \\
Dataset generati & \path{data/pyg}, \path{data/pyg_puzzles_timing} & Ingressi a grafo e ingressi a grafo arricchiti con timing. \\
Artefatti di valutazione & report finali, JSON classici, uscite LLM & Uscite metriche e metadati delle esecuzioni. \\
Analisi finale & \path{generic_info/final_analysis} & Tabelle aggregate destinate al report e inventario. \\
\bottomrule
\end{tabular}
\end{table}


\section{Tre casi d'uso per la riproduzione}
\label{app:reproducibility:workflows}

Il repository supporta tre diversi obiettivi di riproduzione, mostrati nella Tabella~\ref{tab:reproducibility-workflows}.  L'ispezione dei risultati pubblicati richiede soltanto gli artefatti e i manifesti usati dal report.  La valutazione dei modelli pretrained congelati richiede i file dei checkpoint, gli hash esatti, i dataset a grafo e gli entry point di valutazione.  Il riaddestramento usa le configurazioni ufficiali tramite \texttt{train\_official}; è una riesecuzione scientifica con la stessa configurazione prevista, non una garanzia di uguaglianza bit-per-bit tra hardware, CUDA, PyTorch o versioni di PyTorch Geometric.

\begin{table}[H]
\centering
\small
\caption{Tre casi d'uso distinti per la riproduzione.}
\label{tab:reproducibility-workflows}
\begin{tabular}{p{0.20\linewidth}p{0.34\linewidth}p{0.32\linewidth}}
\toprule
Caso d'uso & Artefatti richiesti & Dichiarazione attesa \\
\midrule
Ispezionare i risultati pubblicati & File di analisi finale, report, manifesti & Ricostruire tabelle e provenienza dei risultati riportati. \\
Valutare modelli congelati & Checkpoint, dataset a grafo, manifesti dei benchmark, valutatori & Ricalcolare le metriche per pesi addestrati fissati. \\
Riaddestrare configurazioni ufficiali & Codice sorgente, configurazioni, dati, ambiente di addestramento & Riprodurre la procedura prevista di addestramento, senza garanzia di pesi byte-identici. \\
\bottomrule
\end{tabular}
\end{table}


\section{Interfaccia ufficiale di addestramento}
\label{app:reproducibility:training}

L'entry point ufficiale di addestramento è \texttt{src/cli/training/train\_official.py}.  Accetta \texttt{--model a}, \texttt{a2}, \texttt{a3}, \texttt{b} o \texttt{a4}, oppure un percorso esplicito di configurazione.  La modalità \texttt{--dry-run} ispeziona e valida la configurazione selezionata senza avviare l'addestramento.  Le uscite di riaddestramento sono indirizzate a \texttt{artifacts/retraining/<model>/<timestamp>/}; l'interfaccia rifiuta di scrivere nelle directory canoniche dei checkpoint.  Questa regola di sicurezza protegge l'inventario dei checkpoint congelati da sovrascritture accidentali.

Le dipendenze dati specifiche per modello sono applicate dai loader.  Model B richiede la root PyG con timing invece della root dei grafi senza timing.  Model A4 richiede il checkpoint congelato di Model A3 e la cache delle candidate post-mossa di A4.  Questi requisiti fanno parte della definizione del modello, non sono semplici comodità opzionali.

\section{Riproduzione di checkpoint e valutazioni}
\label{app:reproducibility:evaluation}

Gli hash dei checkpoint congelati sono elencati nell'Appendice~\ref{app:hyperparameters}.  Gli script di valutazione devono verificare fingerprint di dataset e benchmark dove i protocolli ufficiali lo richiedono.  Il valutatore delle composizioni classiche controlla il fingerprint di congelamento \nolinkurl{57ff2b7c725a4ef802817581fac9430f8c6416a6de54dafcb1deb7b471716b60}; il benchmark LLM controlla prompt, parser, digest del modello e identità del dataset prima di accettare un'esecuzione.  Questi controlli rendono l'identità dell'esecuzione più precisa di una sola invocazione da riga di comando.

Gli artefatti di analisi finale aggregano valutazioni completate in tabelle destinate al report.  Sono utili per ispezionare i numeri pubblicati, ma non sostituiscono i file dei checkpoint congelati o i manifesti dei benchmark quando si riproducono le valutazioni sottostanti.

\section{Ambiente e limiti di un clone fresco}
\label{app:reproducibility:environment}

\enlargethispage{2\baselineskip}

Il repository fissa seed, configurazioni ufficiali, manifesti degli split, hash dei checkpoint e fingerprint dei benchmark.  Non garantisce un riaddestramento byte-identico in tutti gli ambienti.  Kernel GPU, versioni delle librerie, comportamento mixed precision, scheduling dei data loader e dettagli implementativi di PyTorch/PyTorch Geometric possono influenzare le traiettorie esatte di ottimizzazione.  Un clone fresco può inoltre richiedere artefatti grandi forniti esternamente o rigenerati, a seconda di quali dataset generati e checkpoint siano inclusi nella distribuzione pubblica.  La dichiarazione di riproducibilità è quindi più forte per l'ispezione degli artefatti e la valutazione dei checkpoint congelati, e più debole per l'equivalenza byte-per-byte del riaddestramento.
